\begin{definition}[6.7.6]
\label{IV.2.s6.7.6-fr}
Sean $k$ un cuerpo y $X$ un $k$-preesquema localmente noetheriano. Se dice
que $X$ es \textup{geométricamente regular en un punto $x$ (resp. posee la
propiedad $(R_n)$ geométrica en un punto $x$, resp. es geométricamente
normal en un punto $x$)} si, para toda extensión finita $k'$ de $k$,
$X'=X_{(k')}=X\otimes_k k'$ es regular \textup{(resp.} posee la propiedad
$(R_n)$\textup{, resp.} es normal) en todo punto $x'$ cuya proyección en
$X$ es $x$. Se dice que $X$ es geométricamente regular \textup{(resp.} posee
la propiedad $(R_n)$ geométrica \textup{(o también} es geométricamente regular
en codimensión $\leqslant n$)\textup{, resp.} es geométricamente normal) si $X$
es geométricamente regular \textup{(resp.} posee la propiedad $(R_n)$
geométrica\textup{, resp.} es geométricamente normal) en todo punto.
\end{definition}

Se dice que un álgebra $A$ sobre $k$ es un \emph{anillo geométricamente regular}
(resp. \emph{geométricamente normal}, resp. \emph{geométricamente reducido}, resp.
\emph{que posee la propiedad $(R_n)$ geométrica}) si $\operatorname{Spec}(A)$
posee la misma propiedad.

Si $X=\operatorname{Spec}(K)$, donde $K$ es una extensión de $k$, es
equivalente decir que $X$ es \emph{geométricamente regular}, o \emph{geométricamente
normal}, o \emph{geométricamente reducido}, o que $K$ es una extensión \emph{separable}
de $k$: esto resulta de (4.6.1) y del hecho de que, si $K$ es una extensión
separable de $k$ y $k'$ una extensión finita de $k$, $K\otimes_k k'$ es un
producto directo de un número finito de cuerpos (Bourbaki, \emph{Alg.}, cap. VIII,
§~7, n\textsuperscript{o}~3, cor. 1 del th. 1).

\oldpage[IV-2]{149}

\begin{proposition}[6.7.7]
\label{IV.2.s6.7.7-fr}
Sean $k$ un cuerpo, $X$ un $k$-preesquema localmente noetheriano y $x$ un
punto de $X$; designemos por $\boldsymbol{Q}(k')$ la relación «$k'$ es una extensión de
$k$, y $\boldsymbol{P}(X_{(k')},x')$ es verdadera para todo punto $x'$ de $X_{(k')}$ del
que $x$ es la proyección en $X$», donde $\boldsymbol{P}(Z,z)$ es una de las propiedades \textup{a),
b), c)} del enunciado \textup{(6.7.4)}. Entonces las propiedades siguientes son
equivalentes:
\begin{enumerate}
\item[\textup{a)}] $\boldsymbol{Q}(k')$ es verdadera para toda extensión finita $k'$ de $k$.
\item[\textup{b)}] $\boldsymbol{Q}(k')$ es verdadera para toda extensión finita radicial $k'$ de $k$.
\item[\textup{c)}] $\boldsymbol{Q}(k')$ es verdadera para toda extensión de tipo finito $k'$ de $k$.
\end{enumerate}
Supongamos además que $X$ es localmente de tipo finito sobre $k$; entonces las
tres propiedades precedentes son también equivalentes a las siguientes:
\begin{enumerate}
\item[\textup{d)}] $\boldsymbol{Q}(k')$ es verdadera para toda extensión $k'$ de $k$.
\item[\textup{e)}] $\boldsymbol{Q}(k')$ es verdadera para \textup{una} extensión perfecta $k'$ de $k$.
\item[\textup{f)}] $\boldsymbol{Q}(k')$ es verdadera para toda extensión $k'=k^{p^{-s}}$, donde $p$ es
el exponente característico de $k$, y $s>0$.
\end{enumerate}
\end{proposition}
